\documentclass[12pt,a4paper]{article}
\usepackage{fontspec}
\setmainfont{cmunrm.otf}[
  BoldFont=cmunbx.otf,
  ItalicFont=cmunti.otf,
  BoldItalicFont=cmunbi.otf
]
\setsansfont{cmunss.otf}[
  BoldFont=cmunsx.otf,
  ItalicFont=cmunsi.otf,
  BoldItalicFont=cmunso.otf
]
\setmonofont{cmuntt.otf}[
  BoldFont=cmuntb.otf,
  ItalicFont=cmunit.otf,
  BoldItalicFont=cmuntx.otf
]
\usepackage[english,russian]{babel}
\usepackage[left=30mm,right=20mm,top=20mm,bottom=20mm]{geometry}
\usepackage{indentfirst}
\usepackage{setspace}
\usepackage{enumitem}
\usepackage{array}
\usepackage{tabularx}
\usepackage{listings}
\usepackage{needspace}
\usepackage{titlesec}
\usepackage[hidelinks]{hyperref}

\setlength{\parindent}{1.25cm}
\setlength{\parskip}{3pt}
\setstretch{1.15}
\setlist{nosep,leftmargin=1.25cm,topsep=5pt}
\titleformat{\section}{\large\bfseries}{}{0pt}{}
\titleformat{\subsection}{\normalsize\bfseries}{}{0pt}{}
\titlespacing*{\section}{0pt}{14pt}{6pt}
\titlespacing*{\subsection}{0pt}{10pt}{4pt}
\newcommand{\code}[1]{\texttt{\detokenize{#1}}}
\newcolumntype{Y}{>{\raggedright\arraybackslash}X}
\renewcommand{\arraystretch}{1.2}
\lstset{
  language=C,
  basicstyle=\ttfamily\footnotesize,
  keywordstyle=\bfseries,
  commentstyle=\itshape,
  numbers=left,
  numberstyle=\tiny,
  numbersep=8pt,
  xleftmargin=14pt,
  breaklines=true,
  breakatwhitespace=false,
  keepspaces=true,
  columns=fullflexible,
  showstringspaces=false,
  tabsize=4,
  aboveskip=7pt,
  belowskip=7pt
}
\hypersetup{pdftitle={Отчёт по лабораторной работе № 2. Вариант 19},
  pdfauthor={Прудников Я. А.}}
\raggedbottom
\begin{document}

\begin{center}
  {\large\bfseries Отчёт по лабораторной работе № 2}\\[4pt]
  по курсу «Операционные системы»\\[4pt]
  Вариант 19
\end{center}

\noindent Студент группы М8О-203БВ-25: Прудников Я. А.

\section*{Тема}
Управление процессами в ОС Linux. Обмен данными между процессами с помощью каналов.

\section*{Цель}
Получить практические навыки создания дочерних процессов, передачи данных через каналы и перенаправления стандартных потоков ввода и вывода.

\section*{Задание}
Написать программу на языке Си, в которой родительский процесс создаёт два дочерних процесса. Родительская и дочерняя программы должны быть отдельными исполняемыми файлами.

Первыми двумя строками пользователь вводит имена файлов для записи результатов первого и второго дочерних процессов. После этого пользователь вводит строки произвольной длины. С вероятностью 80\% родитель отправляет строку в \code{pipe1}, иначе в \code{pipe2}. Дочерние процессы удаляют гласные из полученных строк и записывают результат через стандартный вывод в соответствующие файлы. Необходимо предусмотреть обработку системных ошибок.

\section*{Идея, метод, алгоритм}
Используем две программы: \code{parent} и \code{child}. Родитель запускает два экземпляра \code{child}. У каждого экземпляра свой входной канал и свой выходной файл, заданный пользователем.

\begin{enumerate}
  \item Считываем имена двух файлов и удаляем из них символ перевода строки.
  \item Создаём два неименованных канала.
  \item Создаём два дочерних процесса. Их стандартный ввод перенаправляем на чтение из его канала, а стандартный вывод на запись в его файл.
  \item Запускаем программу \code{child} в каждом дочернем процессе после \code{fork}.
  \item Считываем очередную строку в родительском процессе. Получаем число \texttt{rand() \% 100}. Если оно меньше 80, отправляем строку в первый канал, иначе во второй.
  \item Дочерний процесс читает строку, удаляет из неё гласные и выводит результат.
  \item После окончания ввода закрываем оба конца каналов для записи. Дочерние процессы получают конец потока и завершаются
  \item Дожидаемся завершения обоих процессов, проверяем их коды выхода и освобождаем ресурсы.
\end{enumerate}

\section*{Реализация}
\subsection*{Организация обмена данными}
Стандартный ввод родителя связан с пользовательским вводом. В дочерних процессах \code{stdin} связан с каналом, а \code{stdout} с выходным файлом. Поэтому программа \code{child} работает с обычными потоками и не выбирает канал или файл самостоятельно.

Каналы создаём через \code{pipe2} с флагом \code{O_CLOEXEC}. Это сделано для того, чтобы каналы, унаследованные от родителя файловые дескрипторы при вызове \code{exec} автоматически закрывались. Процессы создаём через \code{fork}, потоки перенаправляем через \code{dup2}, а программу \code{child} запускаем через \code{execvp}. После замены программы исходные дескрипторы каналов закрываются благодаря \code{O_CLOEXEC}. Перенаправленные стандартные потоки остаются открытыми.

Файлы открываем с флагами \code{O_WRONLY | O_CREAT | O_TRUNC}. Если файл уже существует, его прежнее содержимое удаляется. Для результатов используем два разных имени файла. Для ожидания используем \code{waitpid}.

\subsection*{Обработка строк}
В функции \code{read_line} сначала выделяем буфер на 256 байт. Если места не хватает, увеличиваем его в два раза через \code{realloc}. Это позволяет читать строки без заранее заданного ограничения длины, пока хватает памяти. Символ перевода строки сохраняем. Последнюю строку обрабатываем и в случае, если она заканчивается сразу концом потока

В функции \code{remove_vowels} используем два индекса. Первым просматриваем исходную строку, вторым отмечаем место для следующего сохраняемого символа. Если символ не является гласной, записываем его на это место. В конце ставим нулевой символ.

\subsection*{Распределение и запись строк, завершение работы}
Значения от генератора от 0 до 79 направляют строку в первый канал, остальные значения во второй. Такое правило задаёт указанное в условии правило распределения

Функция \code{os_pipe_write} передаёт строку полностью. Если \code{write} записал только часть данных, продолжаем запись оставшейся части. При завершении вызываем \code{cleanup}: функция закрывает все записи в каналы, ожидает процессы, освобождает строки с именами файлов, структуры каналов и процессов

Проверяем результаты выделения памяти, создания каналов и процессов, открытия файлов, перенаправления потоков, записи и ожидания. При обнаружении ошибки возвращаем ненулевой код завершения. Для ряда системных ошибок выводим сообщение через \code{perror}.

\subsection*{Файловая структура программы}
\begin{itemize}
  \item \code{src/parent.c}: чтение пользовательского ввода, выбор канала и завершение работы
  \item \code{src/child.c}: чтение строк, удаление гласных и вывод результата.
  \item \code{src/line_utils.c}: утилита для чтения и работы со строками
  \item \code{src/platform/os_linux.c}: обёртка над os\_api для ОС Linux.
  \item \code{src/include/os_api.h} и \code{src/include/line_utils.h}: объявления функций и типов
\end{itemize}

\section*{Код}
\subsection*{Алгоритм выбора канала в родительской программе (фрагмент \code{parent.c})}
\begin{lstlisting}
int choice = rand() % 100;
if (choice < 80) {
    if (os_pipe_write(pipe1, line, length) < 0) {
        cleanup(pipe1, pipe2, process_1, process_2,
                file1, file2);
        free(line);
        return 1;
    }
    free(line);
} else {
    if (os_pipe_write(pipe2, line, length) < 0) {
        cleanup(pipe1, pipe2, process_1, process_2,
                file1, file2);
        free(line);
        return 1;
    }
    free(line);
}
\end{lstlisting}

\subsection*{Перенаправление потоков в дочернем процессе (фрагмент \code{os_linux.c})}

\begin{lstlisting}
if (close(input->fd[1]) == -1) {
    perror("close");
    _exit(1);
}
if (dup2(input->fd[0], STDIN_FILENO) == -1) {
    perror("dup2");
    _exit(1);
}
if (input->fd[0] != STDIN_FILENO) {
    if (close(input->fd[0]) == -1) {
        perror("close");
        _exit(1);
    }
}
int file_fd = open(output_file, O_WRONLY | O_CREAT | O_TRUNC, 0644);
if (file_fd == -1) {
    perror("open");
    _exit(1);
}
if (dup2(file_fd, STDOUT_FILENO) == -1) {
    perror("dup2");
    _exit(1);
}
if (file_fd != STDOUT_FILENO) {
    if (close(file_fd) == -1) {
        perror("close");
        _exit(1);
    }
}
char *args[] = {(char *)executable, NULL};
execvp(executable, args);
perror("execvp");
_exit(127);
\end{lstlisting}

\subsection*{Дочерняя программа (\code{child.c})}

\begin{lstlisting}
#include "include/line_utils.h"
#include <stdio.h>
#include <stdbool.h>
#include <stdlib.h>

int main() {
    while (true) {
        char *line = NULL;
        size_t length = 0;
        int result = read_line(stdin, &line, &length);
        if (result == 0) {
            free(line);
            break;
        }
        if (result == -1) {
            free(line);
            return 1;
        }
        remove_vowels(line);
        if (fputs(line, stdout) == EOF) {
            free(line);
            return 1;
        }
        free(line);
    }
    return 0;
}
\end{lstlisting}

\subsection*{Удаление гласных (фрагмент \code{line_utils.c})}
\begin{lstlisting}
void remove_vowels(char *line) {
    if (line == NULL) {
        return;
    }
    size_t j = 0;
    for (size_t i = 0; line[i] != '\0'; i++) {
        char ch = line[i];
        int is_vowel = (ch == 'a' || ch == 'e' ||
            ch == 'i' || ch == 'o' || ch == 'u' ||
            ch == 'A' || ch == 'E' || ch == 'I' ||
            ch == 'O' || ch == 'U');
        if (!is_vowel) {
            line[j] = ch;
            j++;
        }
    }
    line[j] = '\0';
}
\end{lstlisting}

\section*{Тестирование}
\begin{center}
\small
\begin{tabularx}{\linewidth}{|Y|Y|}
\hline
Входные данные & Полученный результат \\
\hline
\code{Hello World} & \code{Hll Wrld} \\
\code{Operating Systems} & \code{prtng Sstms} \\
\code{AEIOU aeiou} & Один пробел, затем перевод строки \\
\hline
Пустая строка с переводом строки & Пустая строка с переводом строки \\
\hline
3000 букв \code{a}, затем \code{BC} и перевод строки & \code{BC} и перевод строки \\
\hline
\code{queue} без перевода строки в конце & \code{q} без перевода строки в конце \\
\hline
Пустой входной поток & Пустой вывод \\
\hline
Два имени файлов, затем конец ввода & Код \code{0}; оба файла пустые \\
\hline
Строка из 3000 цифр \code{0} & Одна строка длиной 3000 символов; код завершения \code{0} \\
\hline
Первое имя файла пустое, второе корректное & Код завершения \code{1}; файлы не созданы \\
\hline
Первый файл в несуществующей папке, второй путь корректный & Ошибка открытия файла; код завершения \code{1}; второй файл пустой \\
\hline
\end{tabularx}
\end{center}

\section*{Вывод}
В ходе работы научился связывать отдельные процессы через неименованные каналы и перенаправлять их стандартные потоки. Разобрался, что после \code{fork} нужно настроить файловые дескрипторы, а \code{execvp} запускает другую программу в уже созданном процессе. На примере завершения ввода понял, почему сначала нужно закрывать запись в каналы и только потом ждать дочерние процессы: иначе они могут продолжать ожидать данные.

Также закрепил работу с динамической памятью при чтении длинных строк и обработку строки без дополнительного буфера.

\end{document}
